\section{The Study in One Picture}
\label{sec:overview}

\begin{figure}[t]
\centering
\begin{tikzpicture}[
  node distance=5mm and 7mm,
  box/.style={draw, rounded corners=2pt, align=center, font=\footnotesize, inner sep=4pt, minimum height=9mm},
  env/.style={box, minimum width=23mm},
  arr/.style={-{Stealth[length=2mm]}, thick}]
  \node[box] (src) {One Rust/Rayon crate\\filter, histograms};
  \node[box, right=of src] (rustc) {rustc nightly\\same release profile};
  \node[env, right=14mm of rustc, yshift=11mm, draw=native, text=native] (nat) {native x86-64\\SSE2};
  \node[env, right=14mm of rustc, yshift=-11mm] (wasm) {wasm32\\threads, simd128};
  \node[env, right=of wasm, yshift=7mm, draw=chrome, text=chrome] (chr) {Chrome (V8)};
  \node[env, right=of wasm, yshift=-7mm, draw=firefox, text=firefox] (ffx) {Firefox\\(SpiderMonkey)};
  \draw[arr] (src) -- (rustc);
  \draw[arr] (rustc) -- (nat.west);
  \draw[arr] (rustc) -- (wasm.west);
  \draw[arr] (wasm) -- (chr.west);
  \draw[arr] (wasm) -- (ffx.west);
  \node[box, below=20mm of src] (camp) {Campaign: seeded\\random order};
  \node[box, right=of camp] (run) {Run = fresh process\\warm-up, 10 iterations};
  \node[box, right=of run] (raw) {Raw: CSV + context\\per run, checksum};
  \node[box, right=of raw] (ana) {Analysis: median per run,\\geometric mean, bootstrap};
  \node[box, right=of ana] (res) {10 primary\\contrasts, Holm};
  \draw[arr] (camp) -- (run);
  \draw[arr] (run) -- (raw);
  \draw[arr] (raw) -- (ana);
  \draw[arr] (ana) -- (res);
  \draw[arr, dashed] (wasm.south) |- ($(run.north)+(0,4mm)$) -- (run.north);
\end{tikzpicture}
\caption{The study. Top: one source and one compiler give two artifacts and
three environments. Bottom: a campaign runs every configuration of the three
environments in one random order; each run is a fresh process that records
its times, its checksum and its context; the analysis is fixed before the
data.}
\Description{A flow from one Rust crate through the compiler to a native binary and a WASM module that runs in Chrome and Firefox, and a second flow from campaign to runs to raw data to analysis to contrasts.}
\label{fig:pipeline}
\end{figure}

Figure~\ref{fig:pipeline} is the whole design. One Rust crate holds the
workloads. The same compiler and release profile build a native x86-64
runner and a WASM module; the module runs in Chrome and in Firefox. A
campaign is a list of configurations (environment, workload, input, thread
count, repetition) shuffled once with a recorded seed, so native, Chrome and
Firefox runs are interleaved. Each run is one fresh process or one fresh
browser launch; it measures itself, checks its output against a checksum and
stores its times with the state of the machine. The analysis reads the runs,
never the iterations, as independent observations.

\subsection{Three builds}
\label{sec:builds-history}

The measured code was built three times, and the study was measured in full
twice. Code layout alone moves timings~\cite{mytkowicz2009}, so a rule holds
from the first build on: every campaign runs one tagged build, and no table,
figure or comparison combines runs of two builds.

\begin{description}[leftmargin=1.5em]
\item[v1.] The first build ran the main campaign and the first exploratory
ones. A review of the measured path on 2026-11-09 then found two asymmetries
that our own code had introduced into RQ2. The shared-counter histograms
indexed a table whose length the compiler did not know, so every pixel paid
a bounds check that the local control did not pay; and the final step of a
histogram (the reduce of the local variant, the read-back of the shared
counters) was timed in some variants and not in others.
\item[\code{measured/v2}.] The second build removed both asymmetries
(fixed-size counter tables; every variant times its final step), added
settings the next campaigns needed (passes per timed sample, a grain per
task, a second fork-join probe, the memory size) and nothing that optimizes.
Every campaign was run again on it. An audit of the finished study on
2026-11-20 found no reason to measure again, but listed what v2 could only
declare: the speedup had no sequential reference, the two histograms did not
build their tasks through the same code, the slow one-thread runs had no
established cause, the compiler tier that the campaigns measure was
inferred, and several explanations were hypotheses.
\item[\code{measured/v3}.] The third build turns those declared items into
measurements. It adds sequential versions of the kernels, sends the atomic
histogram through the same task-building code as the local one, and records
the memory size before the measured phase; the harness also records, per
run, the busy time of every measured CPU, and can run a browser with one
compiler only. No change is an optimization, and the code that counts is
the same machine code as in v2 (Section~\ref{sec:codegen}).
\end{description}

The results of v2 were known when v3 was designed, so v3 is a replication
with corrections, not a blind experiment. To keep the choice of the reported
build from following the numbers, the rule was written before any v3 run
existed: v3 is reported if, by 2026-11-27, its main campaign has every
planned run, no checksum conflict, and at least 10 runs left in every
configuration of the primary contrasts after the plan's exclusions;
otherwise v2 is reported and v3 is not used. The criteria read counts of
runs, never a time or a contrast. They were met on 2026-11-22: 5310 of 5310
runs, none failed, missing or excluded, no checksum conflict.

\paragraph{The replication agrees.} All ten primary decisions of v3 are
those of v2: the two contrasts not rejected in v2 are not rejected in v3
(Chrome on the filter, 0.996 in v3 and 0.990 in v2; Firefox at $k=1$, 1.010
and 0.994), and the other eight are rejected in both with effects of the
same sign. Every exploratory campaign of v2 has its twin on v3. This report
gives the v3 numbers only.
